\chapter{Attacchi Denial-of-Service}

\section{Introduzione}

\begin{defbox}[Denial-of-Service (DoS)]
Un attacco Denial-of-Service (DoS) è un tentativo di compromettere la disponibilità di un servizio ostacolandone o bloccandone completamente la fornitura.
L'attacco tenta di esaurire alcune risorse critiche associate al servizio (ad esempio, larghezza di banda o memoria), impedendogli di rispondere alle richieste degli utenti.
\end{defbox}

\begin{warnbox}[Distributed Denial-of-Service (DDoS)]
Gli attacchi DDoS raggiungono un'efficacia nettamente superiore sfruttando molteplici sistemi informatici compromessi come fonti di attacco simultanee. Queste reti sono costituite da computer e dispositivi (noti come \textit{bot} o \textit{zombie}) infettati da malware, che ne permette il controllo a distanza da parte di un attaccante tramite una \textit{botnet}.
\end{warnbox}

\section{TCP SYN Flood}

Un attacco SYN flood è un tipo di DoS che mira a rendere un server indisponibile consumandone tutte le risorse di connessione. Inviando ripetutamente pacchetti di richieste di connessione iniziale (SYN), l'attaccante esaurisce le porte o la tabella di stato del server, obbligandolo a rispondere lentamente o a non rispondere affatto al traffico legittimo.

\subsection{Funzionamento del TCP SYN Flood}

\subsubsection{Funzionamento Base di una Connessione TCP}

Gli attacchi SYN flood sfruttano il processo di \textit{three-way handshake} di una connessione TCP.
\begin{enumerate}
    \item Il client invia un pacchetto \texttt{SYN} al server per avviare la connessione.
    \item Il server risponde con un pacchetto \texttt{SYN/ACK} confermando la ricezione e accettando la comunicazione.
    \item Il client restituisce un pacchetto \texttt{ACK} per confermare a sua volta.
\end{enumerate}

\subsubsection{Provocare un Denial-of-Service}

Per provocare il DoS, l'attaccante sfrutta il fatto che dopo aver ricevuto il pacchetto \texttt{SYN} iniziale, il server alloca risorse per la connessione (che entra nello stato \textit{semi-aperta} o \textit{SYN-RECEIVED}), risponde con \texttt{SYN/ACK} e attende l'ultimo \texttt{ACK}.

\subsection{Come si Verifica l'Attacco}
\begin{enumerate}
    \item L'attaccante invia un elevato volume di pacchetti \texttt{SYN} al server, spesso falsificando (\textit{spoofing}) l'indirizzo IP di origine.
    \item Il server risponde a ciascuna richiesta con \texttt{SYN/ACK} e alloca una porta/struttura dati, in attesa del completamento.
    \item Poiché l'IP è falsificato (o l'attaccante semplicemente ignora il \texttt{SYN/ACK}), l'atteso pacchetto \texttt{ACK} non arriverà mai. Il server mantiene la connessione aperta fino al timeout. Saturate le risorse disponibili, il server smette di accettare nuove connessioni.
\end{enumerate}

\subsection{Tipologie di Attacco SYN Flood}
\begin{description}
    \item[Attacco diretto] L'indirizzo IP dell'attaccante \textbf{non} è spoofato. È raro, poiché l'attaccante è vulnerabile e la mitigazione è banale (blocco dell'IP). Spesso usato solo se l'attaccante opera tramite botnet, dove l'IP esposto è quello del dispositivo infetto.
    \item[Attacco di spoofing] L'attaccante falsifica l'indirizzo IP nei pacchetti \texttt{SYN}. Rendere più difficile l'identificazione. Risalire all'origine richiede complesse indagini con i provider di servizi Internet (ISP).
    \item[Attacco distribuito (DDoS)] L'attacco è sferrato tramite una botnet e \textbf{ogni} bot falsifica ulteriormente gli IP sorgente. Rintracciare l'origine o mitigare l'attacco tramite semplici filtri IP diventa quasi impossibile.
\end{description}

\subsection{Mitigazione degli Attacchi SYN Flood}
Le moderne infrastrutture mitigate in cloud (es. Cloudflare) si interpongono tra il server vittima e l'attaccante. Gestiscono il processo di \textit{three-way handshake} esternamente. La connessione viene passata al server reale \textbf{solo se e quando} il client invia il pacchetto \texttt{ACK} finale. Questa strategia sposta il carico del SYN flood sull'infrastruttura globale di mitigazione, risparmiando il server.

\section{ICMP Flood}

\begin{defbox}[ICMP Flood / Ping Flood]
Un ICMP flood è un attacco DoS volumetrico in cui l'attaccante sommerge un dispositivo target con pacchetti \texttt{ICMP Echo-Request} (Ping), saturandone la banda o la capacità di elaborazione.
\end{defbox}

\subsection{Funzionamento dell'ICMP Flood}
Il protocollo ICMP (Internet Control Message Protocol) è utilizzato dai dispositivi per comunicazioni diagnostiche (come i tool \texttt{ping} o \texttt{traceroute}). L'attacco mira a costringere il bersaglio a elaborare e rispondere a un volume insostenibile di richieste.

\begin{enumerate}
    \item L'attaccante invia un flusso massiccio di pacchetti \texttt{ICMP Echo-Request} verso l'IP vittima.
    \item Il server bersaglio, per specifiche di protocollo, è tenuto a inviare un pacchetto \texttt{ICMP Echo-Reply} per ciascuna richiesta ricevuta.
    \item Questo doppio lavoro (ricezione massiccia ed emissione massiccia) satura sia le risorse CPU del bersaglio sia la larghezza di banda del suo collegamento di rete.
\end{enumerate}

\subsection{Mitigare l'ICMP Flood}
La mitigazione più semplice consiste nel limitare (rate-limiting) o disabilitare del tutto la risposta ai pacchetti ICMP sul router di frontiera o sul firewall del dispositivo vittima, eliminando l'onere dell'elaborazione della risposta.

\section{UDP Flood}

Un UDP flood è un attacco DoS volumetrico in cui un gran numero di pacchetti UDP (User Datagram Protocol) viene inviato a porte casuali di un server bersaglio, con l'obiettivo di esaurirne la capacità elaborativa o saturare la tabella di stato del firewall.

\subsection{Funzionamento dell'UDP Flood}
UDP è un protocollo \textit{connectionless}, non richiede handshake. Quando un server riceve un datagramma UDP, esegue due passaggi:
\begin{enumerate}
    \item Controlla se c'è un'applicazione (servizio) in ascolto sulla porta di destinazione.
    \item Se nessuna applicazione è in ascolto, il sistema operativo risponde inviando al mittente un pacchetto \texttt{ICMP Destination Unreachable}.
\end{enumerate}

\begin{exbox}[Il Paragone dell'UDP Flood]
Si può paragonare a una receptionist d'albergo che riceve telefonate continue per camere diverse. Per ogni chiamata, deve controllare l'elenco (la tabella delle porte), verificare che l'ospite non ci sia o non risponda, e poi informare il chiamante (il pacchetto ICMP). Se i centralini si illuminano contemporaneamente di richieste fittizie, la receptionist è subito sopraffatta e non può gestire i veri clienti.
\end{exbox}

\subsection{Mitigare l'UDP Flood}
I sistemi operativi moderni limitano per default la velocità di generazione dei pacchetti ICMP. Tuttavia, se il volume di traffico UDP è enorme, il \textit{firewall stateful} a monte del server andrà comunque in saturazione cercando di tracciare tutte queste "finte connessioni". La mitigazione efficace richiede filtraggi a monte (presso l'ISP o servizi anti-DDoS cloud) per assorbire il volume prima che raggiunga il target.

\section{HTTP Flood}

\begin{warnbox}[HTTP Flood (Livello 7)]
È un attacco DoS applicativo (Livello 7 OSI) progettato per sommergere un server Web con richieste HTTP apparentemente legittime. È estremamente insidioso perché è difficile distinguere le richieste malevole da un improvviso picco di traffico reale (effetto \textit{flash crowd}).
\end{warnbox}

\subsection{Tipologie di Attacco HTTP Flood}
\begin{description}
    \item[HTTP GET] Botnet coordinate inviano continue richieste per scaricare risorse pesanti (immagini, PDF, o pagine dinamiche) costringendo il server web a consumare CPU e I/O disco per recuperarle e inviarle.
    \item[HTTP POST] Spesso diretto verso moduli di login o di ricerca. Il server deve gestire i parametri e spesso interrogare un database di backend. L'elaborazione del DB è computazionalmente costosa, quindi con poche richieste POST ben assestate l'attaccante può paralizzare l'intero stack applicativo (esaurimento risorse DB).
\end{description}

\subsection{Mitigare un HTTP Flood}
Implementare una "challenge" (es. CAPTCHA, controlli JavaScript invisibili o analisi comportamentale) prima di consentire l'accesso al server backend, per verificare se il client è un browser umano o un bot.

\section{DNS Amplification Attack}

Gli attacchi di amplificazione sono una combinazione sofisticata di \textit{reflection} (riflessione) e \textit{amplification} (amplificazione), solitamente basati su protocolli UDP.

\subsection{Concetti Chiave}
\begin{description}
    \item[Reflection Attack] L'attaccante invia richieste a server legittimi terzi (riflettori) \textbf{falsificando l'IP sorgente} per farlo coincidere con quello della vittima. Il server terzo elaborerà la richiesta e invierà la risposta alla vittima.
    \item[Amplification Attack] Sfrutta servizi in cui una richiesta piccola (\textit{trigger packet}) genera una risposta molto più grande (fattore di amplificazione). L'attaccante moltiplica l'efficacia della sua banda.
\end{description}

\subsection{Funzionamento del DNS Amplification Attack}
Sfrutta i server DNS (Open Resolvers) per amplificare enormemente il traffico contro un bersaglio.

\begin{center}
\begin{tikzpicture}[node distance=2cm, auto, thick, scale=0.9, transform shape]
\node (att) [draw, rounded corners, fill=red!20, text width=2.5cm, align=center, minimum height=1cm] {Attaccante};
\node (dns) [draw, rounded corners, fill=green!20, text width=2.5cm, align=center, right=3cm of att, minimum height=1cm] {Server DNS\\(Riflettore)};
\node (vic) [draw, rounded corners, fill=blue!10, text width=2.5cm, align=center, below=2cm of dns, minimum height=1cm] {Vittima};

\draw[->, >=Stealth, red, thick] (att) -- node[above, text width=2.5cm, align=center, font=\footnotesize] {1. Query piccola\\IP Src = Vittima} (dns);
\draw[->, >=Stealth, red, line width=2pt] (dns) -- node[right, text width=2.8cm, align=left, font=\footnotesize] {2. Risposta enorme\\(Amplificata)} (vic);
\end{tikzpicture}
\end{center}

\begin{enumerate}
    \item L'attaccante invia (tramite botnet) query UDP falsificando l'IP mittente con l'IP della vittima verso numerosi server DNS pubblici (Open Resolvers).
    \item La richiesta chiede record molto estesi (es. parametro \texttt{ANY}, firme DNSSEC o ampi record TXT/IPv6).
    \item I server DNS, comportandosi correttamente secondo il protocollo, inviano le risposte massive all'indirizzo IP che credono essere il richiedente (la vittima).
    \item La rete della vittima viene invasa da un volume di traffico che è decine di volte superiore a quello generato originariamente dall'attaccante (fattore di amplificazione fino a 50x o più).
\end{enumerate}

\subsection{Mitigare il DNS Amplification Attack}
\begin{itemize}
    \item \textbf{Ingress Filtering (BCP 38)}: I provider di rete dovrebbero scartare i pacchetti in uscita che hanno IP sorgente non appartenenti alla propria sottorete, impedendo di fatto l'IP spoofing.
    \item \textbf{Sicurezza dei Server DNS}: Disabilitare la ricorsione pubblica sui server DNS (\textit{Open Resolvers}). Un DNS dovrebbe fornire ricorsione solo agli utenti della propria rete locale, non a tutta Internet.
\end{itemize}
